% =====================================================================
\section{Evaluation and Testing}
\label{sec:evaluation}
% =====================================================================

\subsection{Metrics}
The proposal specifies Word Error Rate and Character Error Rate as the
evaluation metrics \spec. Both are based on the minimum edit distance between
the recognised text (hypothesis) and the reference transcript. For a reference of
$N$ words, let $S$, $D$ and $I$ be the numbers of word substitutions, deletions
and insertions in the minimum-cost alignment. Then
\begin{equation}
\mathrm{WER} = \frac{S + D + I}{N}\times 100\%,
\qquad
\mathrm{CER} = \frac{S_c + D_c + I_c}{N_c}\times 100\%,
\label{eq:wer}
\end{equation}
where the subscript $c$ denotes the same quantities counted over characters
(including spaces). Corpus-level values are computed by summing errors and
reference lengths over the whole test set before dividing, not by averaging
per-utterance rates; this is the convention of the \texttt{jiwer} library used
for the computation. WER can exceed 100\% when there are many insertions.
CER is the more sensitive metric for GRID's single-letter slot, where a single
wrong character makes a whole word wrong.

\subsection{Evaluation protocol}
\begin{enumerate}
  \item Load the best checkpoint of each model (selected on validation data
        only).
  \item For each condition $\rho \in \{\text{clean}, 10, 5, 0, -5\}$\,dB, decode
        every test utterance of the four held-out speakers with beam search,
        using the fixed pre-mixed noisy audio.
  \item Compute corpus-level WER and CER per condition, per speaker and in
        total, and record the mean fusion gate per utterance.
  \item Repeat for every model in \cref{tab:experiments} on identical inputs.
  \item Where three seeds were trained, report mean and standard deviation.
\end{enumerate}

Performance is reported separately for each SNR, because averaging across noise
levels would hide the effect the project is designed to measure: how quickly
accuracy degrades as noise increases and where the visual stream begins to
help. The main result will be a WER-versus-SNR plot with one line per model and a
matching table, as sketched in \cref{tab:results-template}.

\begin{table}[H]
\centering
\caption{Template for reporting test WER (\%). Cells are intentionally empty: no
results have been obtained yet.}
\label{tab:results-template}
\small
\begin{tabular}{@{}lccccc@{}}
\toprule
\textbf{Model} & \textbf{Clean} & \textbf{10\,dB} & \textbf{5\,dB} & \textbf{0\,dB} & \textbf{$-5$\,dB} \\
\midrule
B1 Audio-only        & -- & -- & -- & -- & -- \\
B2 Visual-only       & -- & -- & -- & -- & -- \\
B3 Concatenation     & -- & -- & -- & -- & -- \\
P  Proposed          & -- & -- & -- & -- & -- \\
\bottomrule
\end{tabular}
\end{table}

\subsection{Comparisons and their interpretation}
\begin{itemize}
  \item \textbf{Overall quality.} The proposal sets an expected target of WER below
        10\% for the fused model under clean conditions \spec. This is a project
        target, not a guaranteed outcome; the evaluation will report whether it
        is met.
  \item \textbf{Degradation with noise.} The slope of WER against SNR for each model
        shows how robust it is. The proposal expects a clear advantage over the
        audio-only baseline at 0 and $-5$\,dB \spec.
  \item \textbf{Versus audio-only (P vs.\ B1).} The primary test of the project's
        hypothesis.
  \item \textbf{Versus visual-only (P vs.\ B2).} At $-5$\,dB the fused model should be
        no worse than the visual-only model, since in the limit it can ignore the
        audio; if it is worse, fusion is being misled by noisy audio.
  \item \textbf{Versus simple fusion (P vs.\ B3).} Isolates the benefit of attention
        and modality dropout over naive combination.
  \item \textbf{Speaker independence.} All numbers are on unseen speakers.
        Per-speaker WER is also reported, because a mean can hide one speaker
        on whom the system fails.
  \item \textbf{Statistical caution.} With four test speakers, differences of a
        fraction of a percentage point may not be meaningful. Where seeds allow,
        differences are reported with their seed-to-seed spread, and a paired
        bootstrap over test utterances is used to give a confidence interval
        for the most important comparison (P vs.\ B1 at each SNR) \rec.
\end{itemize}

\paragraph{If the fused model is worse than audio-only.} This can happen,
particularly on clean audio, where video adds little information but adds
parameters and noise. The team will check, in order: (1) whether the visual-only
model learned anything (if B2 is near chance, the visual pipeline is broken);
(2) whether the gate values vary with SNR or are constant; (3) whether training
loss is lower but validation loss higher than B1, indicating overfitting to
training speakers' faces; (4) whether modality dropout rates are too high,
starving the audio stream; and (5) whether the visual stream is misaligned in
time (\cref{sec:testing}). A fused model that is slightly worse when clean but
better at low SNR is still consistent with the project's goal and will be
reported as such.

\subsection{Ablation studies}
Ablations X1--X3 (\cref{tab:experiments}) each remove one advanced component
while keeping everything else fixed: modality dropout, attention fusion and the
3D convolutional front-end. They are evaluated with the same per-SNR protocol.
Two additional test-time ablations use the proposed model without retraining:
replacing the visual encoding by $\mathbf{e}_{\text{miss}}$ (audio-only operation)
and replacing the audio encoding by $\mathbf{e}_{\text{miss}}$ (visual-only
operation). They show how much each stream contributes inside the fused model.

\subsection{Latency and real-time feasibility}
For the demonstration, the team will measure on the demonstration machine:
\begin{itemize}
  \item per-frame Face Mesh and crop time;
  \item model forward time and beam-search time for one window;
  \item end-to-end caption latency, from the end of a spoken phrase to the
        appearance of its caption;
  \item the real-time factor $\mathrm{RTF} = t_{\text{processing}} / t_{\text{audio}}$.
\end{itemize}
An RTF below 1 is necessary for live operation; the target for the
demonstration is an RTF of at most 0.5 and caption latency of about 3--4\,s
\rec\ (\cref{sec:application}). Timings are reported as the median and 95th
percentile over at least 100 windows, after a warm-up period.

\subsection{Software testing}
\label{sec:testing}
Tests are written with \texttt{pytest} and run before every merge to the main
branch. Small fixture files (two GRID utterances and one MUSAN noise file) are kept
in the repository so that tests run without the full dataset.

\begin{table}[H]
\centering
\caption{Representative automated tests and acceptance criteria.}
\label{tab:tests}
\small
\begin{tabularx}{\textwidth}{@{}p{1.8cm}p{3.6cm}X@{}}
\toprule
\textbf{Level} & \textbf{Test} & \textbf{Assertion / acceptance criterion} \\
\midrule
Unit & Tokeniser round trip & \texttt{decode(encode(s)) == s} for all GRID word strings; blank never produced by \texttt{encode} \\
Unit & CTC collapse & \texttt{collapse([b,b,0,i,0,n,n]) == "bin"}; \texttt{collapse([e,0,e]) == "ee"} \\
Unit & Noise mixing & measured SNR within 0.01\,dB of target for all five conditions; output length equals input length; peak $\le 1$ \\
Unit & Mouth crop & output shape $(96,96)$, dtype \texttt{uint8}; failed detection yields interpolated box and invalid flag \\
Unit & Feature extraction & 3\,s waveform gives \shape{300,80}; matches librosa within tolerance \\
Unit & Encoder shapes & visual and audio encoders map example inputs to \shape{B,75,512} \\
Unit & Metrics & WER(\texttt{"bin blue at f two now"}, \texttt{"bin blue at s two now"}) $= 1/6$ \\
Unit & Beam search & with $W=1$, beam search equals greedy decoding on random posteriors \\
Unit & Splits & train, validation and test speaker sets are pairwise disjoint \\
Integration & Data loader & one batch has consistent lengths; every item satisfies \cref{eq:ctc-cond} \\
Integration & Overfit check & each model reaches CTC loss $< 0.05$ and CER $=0$ on 8 utterances within 300 steps \\
Integration & Train/infer parity & live preprocessing of a GRID file equals cached features (max error $<10^{-5}$) \\
Integration & Synchronisation & shifting video by 5 frames relative to audio increases validation loss of the trained model (sanity check of alignment) \\
End-to-end & Inference script & \texttt{transcribe.py} on a test clip returns a non-empty lower-case string within the latency budget \\
End-to-end & Demo smoke test & demo starts, processes a recorded clip and displays a caption without error \\
\bottomrule
\end{tabularx}
\end{table}
